%----------
%--- DISCUSSÃO

\chapter{DISCUSSÃO}
\thispagestyle{myheadings}

Este capítulo discute os principais ganhos, desafios e limitações decorrentes da reestruturação
do processo de ETL do BDMalhas. A análise fundamenta-se nos resultados apresentados no
capítulo anterior, buscando relacionar os resultados obtidos às características da arquitetura
proposta e identificar os fatores que podem influenciar sua adoção e evolução.

De modo geral, os resultados indicam que a reestruturação proporcionou melhorias relevantes
na estrutura dos \textit{pipelines} e no armazenamento dos dados. Em relação à manutenibilidade,
a redução do número de nós, dependências e consultas SQL evidencia uma simplificação
estrutural dos fluxos de ETL. Ademais, a utilização do Apache Airflow, associada à parametrização das
consultas e à centralização de parte da lógica de processamento em código, também contribuiu
para reduzir a quantidade de componentes necessários à execução dos fluxos.

Devemos destacara que, embora essas métricas não representem uma medida direta da manutenibilidade, 
elas fornecem indicadores objetivos de redução da complexidade estrutural dos \textit{workflows}. 
A menor quantidade de nós, dependências e consultas pode favorecer a compreensão dos fluxos e reduzir
o número de elementos que precisam ser modificados durante futuras alterações. Além disso,
a utilização de consultas parametrizadas e o reaproveitamento de lógica diminuem a duplicação
de código, contribuindo para uma estrutura mais modular em comparação à implementação
original em SSIS.

Em relação ao desempenho, os resultados obtidos indicam um comportamento favorável da nova
solução nas condições experimentais avaliadas. Para as cargas de CTE e NFE, por exemplo,
foram observadas reduções nos tempos totais de execução e aumento das taxas de processamento.
Entretanto, esses resultados não permitem afirmar uma superioridade absoluta da nova solução em 
relação ao sistema legado, uma vez que as execuções ocorreram em ambientes distintos e processaram
volumes diferentes de dados.

Essa limitação é particularmente relevante nas cargas de EFD e NFCE, nas quais a diferença
entre os volumes processados pelas duas implementações foi bem expressiva.  

(completar) Dessa forma, embora o tempo observado na nova solução tenha sido
consideravelmente menor, essa diferença não pode ser utilizada isoladamente como evidência
de superioridade de desempenho.

Desse modo, os resultados devem ser interpretados como evidências do comportamento da solução no
ambiente experimental utilizado, sendo necessária uma avaliação com volumes equivalentes e 
infraestrutura comparável para uma conclusão mais abrangente.

Outro aspecto relevante é que a nova arquitetura reúne diferentes tecnologias, como Apache
Airflow, Apache Spark e ClickHouse. Os ganhos
identificados devem ser compreendidos como decorrentes da combinação dos componentes da
arquitetura proposta, juntamente com as alterações realizadas nos fluxos de processamento e
no modelo de dados utilizado como fonte.

No que se refere ao armazenamento, os resultados obtidos nas cinco tabelas analisadas
demonstraram uma redução expressiva do espaço ocupado após a migração dos dados para o
ClickHouse. As reduções observadas ficaram entre 78,59\% e 89,81\%, indicando um
comportamento consistente de redução do espaço para os diferentes volumes avaliados.
Esse resultado está relacionado às características do armazenamento colunar e aos mecanismos
de compressão utilizados pelo ClickHouse, que permitem agrupar valores de uma mesma coluna
e explorar padrões e repetições presentes nos dados. Contudo, os resultados são limitados às
cinco tabelas selecionadas e, portanto, não permitem generalizar diretamente a mesma taxa de
redução para todo o banco de dados.

Além dos resultados quantitativos, a reestruturação também apresentou desafios relacionados
à migração da solução baseada em SSIS para uma arquitetura composta por Airflow e Spark.
Embora o Airflow ofereça maior flexibilidade para a implementação programática dos fluxos,
determinadas funcionalidades disponíveis por meio de componentes do SSIS precisaram ser
reimplementadas utilizando código Python. Essa mudança aumenta a flexibilidade da solução, 
mas também transfere parte da responsabilidade de implementação e manutenção para o código
desenvolvido pela equipe.

Outro desafio esteve relacionado à utilização do \textit{Data Lake} como nova fonte de dados.
A diferença entre o esquema disponível no \textit{Data Lake} e as estruturas existentes nas
bases de produção exigiu adaptações nas consultas de extração. A ausência da coluna
\textit{sqcontrib}, utilizada originalmente na carga de EFD, é um exemplo dessa limitação.
Nesse caso, foi necessário utilizar o CNPJ como critério alternativo de filtragem, o que
impediu uma correspondência exata entre os dados processados pelas duas implementações.

As limitações do ambiente experimental também restringem a generalização dos resultados.
A nova solução foi executada sobre uma réplica do \textit{Data Lake} contendo um volume de
dados inferior ao disponível nas bases de produção. Além disso, as execuções das duas
implementações ocorreram em infraestruturas distintas e a solução legada foi executada apenas
uma vez. Consequentemente, os resultados de desempenho representam as condições observadas
nos experimentos realizados e não uma avaliação estatística da variabilidade das execuções.

Diante dessas limitações, a evolução da solução depende de futuras iterações com a equipe
responsável pelo \textit{Data Lake}, especialmente para completar o esquema de dados e
adequar as consultas de extração às estruturas existentes nas bases de produção. A
disponibilização de um conjunto de dados mais completo também permitirá realizar experimentos
com volumes mais próximos dos encontrados no ambiente produtivo, tornando futuras comparações
de desempenho mais representativas.

De forma geral, os resultados indicam que a reestruturação trouxe benefícios principalmente
na simplificação estrutural dos fluxos e na redução do espaço de armazenamento. Os resultados
de desempenho também foram favoráveis em determinados cenários, especialmente para CTE e NFE,
mas devem ser interpretados considerando as diferenças entre os ambientes e os volumes
processados. Assim, a arquitetura proposta demonstra potencial para substituir e evoluir o
processo de ETL existente, embora sua avaliação definitiva em ambiente produtivo dependa da
conclusão das adequações do \textit{Data Lake} e da realização de novos experimentos sob
condições mais equivalentes.

